\scanpage{82}
In the ``arithmetic'' case, one therefore sees that the extension (13) is split if and only if $\gamma_i=0$ that is if one can find, globally on $S$, a uniformizing parameter on $X$ along $\unsure{g|g_i}(S)$; and this condition (always satisfied locally\inserted{(\unsure{Zariski})} on $S$) is in any case always sufficient (arithmetic case or not !)

What is more, among the splittings, there are those coming from the choice of a \deleted{section}\inserted{basis} of $J_i/J_i^2$ which are preferred. The indeterminacy of the choice of such a basis is in $G_m(S)=H^0(S,G_m)$ --- that of the choice of a section of (13) is in
\[
(20)\quad H^1(S,T)\simeq\varprojlim H^1(S,\mu_n)
\]
Here one has Kummer exact sequences
\[
\deleted{\unsure{(21)|(22)}}\quad 0-\deleted{\unsure{H^0(S)_n}}G_m(S)_n\longrightarrow H^1(S,\mu_n)\longrightarrow{}_n\operatorname{Pic}(S)\longrightarrow1
\]
whence by passing to the limit
\[
(21)\quad 0-G_m(S)^\wedge\longrightarrow H^1(S,T)\longrightarrow\varprojlim_n{}_n\operatorname{Pic}S\unsure{\longrightarrow?|\longrightarrow0}
\]
In the ``arithmetic case'' the third term is zero, one therefore finds
\[
\deleted{G_m}\quad H^1(S,T)\simeq G_m(S)^\wedge.
\]
\margin{81}
